\chapter{The IPA Scheme over the Pasta Curves}

The preceding chapters develop the IPA polynomial commitment generically: the
scheme \(\textsf{ipa}\) of \cref{def:ipa_scheme} and its two security theorems are
stated for an arbitrary additive commitment group \(\GG\) of prime order with
scalar field \(\Fr\). This chapter instantiates that generic development at the two
Pasta curves, Pallas and Vesta, taking \(\GG\) to be a Pasta curve's group of
rational points and \(\Fr\) to be the corresponding scalar field, both supplied by
CompElliptic. The result is a concrete \texttt{Commitment.Scheme} on each curve
together with the specialized correctness and binding statements; nothing new is
proved here beyond transporting the generic results across the instantiation.

The Pallas and Vesta curves form a \(2\)-cycle: both have the short-Weierstrass
shape \(y^{2} = x^{3} + 5\), the base field of one is the scalar field of the
other, and each curve's group of points has prime order equal to the cardinality of
the partner field. Concretely, the Pallas curve is defined over the Pallas base
field, and its group of rational points has order \(q\), the cardinality of the
Pallas scalar field; the Vesta curve is defined over that scalar field \(\Fp[q]\),
and its group has order \(p\), the cardinality of the Pallas base field. Each
instantiation pairs a curve's point group with the partner curve's base field as
the scalar field, exactly as the generic scheme requires.

\section{CompElliptic Pasta dependency anchors}

The following four blocks are dependency anchors for the CompElliptic vendor
library: each states an external declaration in the project's notation and points
\lean{} at the real CompElliptic target, marked \mathlibok so the dependency graph
treats it as supplied. None is a proof obligation of this project — the \lean{}
target is the citation.

\begin{definition}[Pallas curve group]
  \label{def:pallas_group}
  \lean{CompElliptic.Curves.Pasta.Pallas.curve}
  \mathlibok
  \textit{Provided by CompElliptic.}
  The Pallas curve is the short-Weierstrass elliptic curve
  \(y^{2} = x^{3} + 5\) (coefficients \(A = 0\), \(B = 5\)) over the Pallas base
  field. Its discriminant is a unit, so the curve is elliptic, and its set of
  rational points forms a finite abelian group of prime order \(q\). This point
  group is the commitment group \(\GG\) of the Pallas instantiation of the IPA
  scheme.
\end{definition}

\begin{definition}[Pallas scalar field]
  \label{def:pallas_scalar_field}
  \lean{CompElliptic.Fields.Pasta.PallasScalarField}
  \mathlibok
  \textit{Provided by CompElliptic.}
  The Pallas scalar field is the prime field \(\Zmod{q}\), where \(q\) is the
  (prime) order of the Pallas point group of \cref{def:pallas_group}. It carries a
  machine-checked Pratt/Lucas primality certificate for \(q\), and is the scalar
  field \(\Fr\) over which polynomials are committed and folding challenges are
  drawn in the Pallas instantiation.
\end{definition}

\begin{definition}[Vesta curve group]
  \label{def:vesta_group}
  \lean{CompElliptic.Curves.Pasta.Vesta.curve}
  \mathlibok
  \textit{Provided by CompElliptic.}
  The Vesta curve is the short-Weierstrass elliptic curve \(y^{2} = x^{3} + 5\)
  (coefficients \(A = 0\), \(B = 5\)) over the Vesta base field, which coincides
  with the Pallas scalar field of \cref{def:pallas_scalar_field}. Its rational
  points form a finite abelian group of prime order \(p\), equal to the cardinality
  of the Pallas base field. This point group is the commitment group \(\GG\) of the
  Vesta instantiation.
\end{definition}

\begin{definition}[Vesta scalar field]
  \label{def:vesta_scalar_field}
  \lean{CompElliptic.Fields.Pasta.VestaScalarField}
  \mathlibok
  \textit{Provided by CompElliptic.}
  The Vesta scalar field is the prime field \(\Zmod{p}\), where \(p\) is the order
  of the Vesta point group of \cref{def:vesta_group}; it coincides with the Pallas
  base field. It is the scalar field \(\Fr\) of the Vesta instantiation.
\end{definition}

\section{The Pasta instances}

\begin{definition}[IPA scheme over Pallas]
  \label{def:ipa_pallas}
  \lean{Kimchi.Commitment.IPA.ipaPallas}
  \uses{def:ipa_scheme, def:pallas_group, def:pallas_scalar_field}
  The IPA commitment scheme over the Pallas curve is the generic scheme
  \(\textsf{ipa}\) of \cref{def:ipa_scheme} with the commitment group \(\GG\) taken
  to be the Pallas point group of \cref{def:pallas_group} and the scalar field
  \(\Fr\) taken to be the Pallas scalar field of \cref{def:pallas_scalar_field}.
  All of the abstract data slots of \cref{def:ipa_scheme} — the structured
  reference string of generators and blinding base, the Pedersen commitment, and
  the logarithmic-round opening protocol — are inherited unchanged; only the
  carrier group and its scalars are fixed. Concretely, a polynomial of degree
  \(< d\) over the Pallas scalar field is committed by a Pedersen vector commitment
  whose values lie in the Pallas point group.
\end{definition}

\begin{definition}[IPA scheme over Vesta]
  \label{def:ipa_vesta}
  \lean{Kimchi.Commitment.IPA.ipaVesta}
  \uses{def:ipa_scheme, def:vesta_group, def:vesta_scalar_field}
  The IPA commitment scheme over the Vesta curve is the generic scheme
  \(\textsf{ipa}\) of \cref{def:ipa_scheme} with the commitment group \(\GG\) taken
  to be the Vesta point group of \cref{def:vesta_group} and the scalar field
  \(\Fr\) taken to be the Vesta scalar field of \cref{def:vesta_scalar_field}. As
  for Pallas, every component of the generic scheme is inherited unchanged; the
  construction differs only in the choice of carrier curve. Because the Pasta
  curves form a cycle, the Pallas and Vesta instantiations are the two halves used
  together by the recursive proof system, but each is, on its own, an instance of
  the same generic IPA scheme.
\end{definition}

\section{Security at the Pasta curves}

\begin{theorem}[perfect correctness over Pallas]
  \label{thm:ipa_pallas_correctness}
  \lean{Kimchi.Commitment.IPA.ipaPallas_correctness}
  \uses{thm:ipa_correctness, def:ipa_pallas, def:commitment_perfect_correctness}
  The Pallas IPA scheme of \cref{def:ipa_pallas} satisfies perfect correctness
  (\cref{def:commitment_perfect_correctness}): for every polynomial of degree
  \(< d\) over the Pallas scalar field and every evaluation point, the honest
  commit-then-open interaction makes the verifier accept with probability \(1\).
  \begin{proof}
    \uses{thm:ipa_correctness}
    The generic perfect-correctness theorem \cref{thm:ipa_correctness} holds for
    the scheme \(\textsf{ipa}\) over any prime-order commitment group with its
    scalar field. The Pallas instantiation of \cref{def:ipa_pallas} is exactly
    \(\textsf{ipa}\) with \(\GG\) the Pallas point group
    (\cref{def:pallas_group}) and \(\Fr\) the Pallas scalar field
    (\cref{def:pallas_scalar_field}); both hypotheses of the generic theorem —
    that \(\GG\) is a commutative group acted on by the field \(\Fr\) and that
    \(\Fr\) is a field — are discharged by the CompElliptic instances for the
    Pasta curve. Specializing \cref{thm:ipa_correctness} at these choices yields
    the claim.
  \end{proof}
\end{theorem}

\begin{theorem}[evaluation binding over Pallas]
  \label{thm:ipa_pallas_binding}
  \lean{Kimchi.Commitment.IPA.ipaPallas_binding}
  \uses{thm:ipa_binding, def:ipa_pallas, def:commitment_binding}
  The Pallas IPA scheme of \cref{def:ipa_pallas} satisfies evaluation binding
  (\cref{def:commitment_binding}): no efficient adversary can open a single Pallas
  commitment to two different evaluations at the same point, except with the
  negligible error of the generic bound.
  \begin{proof}
    \uses{thm:ipa_binding}
    The generic binding theorem \cref{thm:ipa_binding} establishes evaluation
    binding for \(\textsf{ipa}\) over any prime-order commitment group with its
    scalar field, the binding error being controlled by the hardness of the
    discrete-logarithm relation among the SRS elements. The Pallas scheme is the
    instantiation of \(\textsf{ipa}\) at the Pallas point group
    (\cref{def:pallas_group}) and Pallas scalar field
    (\cref{def:pallas_scalar_field}), whose prime order is certified by
    CompElliptic; specializing \cref{thm:ipa_binding} at this group and field
    gives the result.
  \end{proof}
\end{theorem}

\begin{theorem}[perfect correctness over Vesta]
  \label{thm:ipa_vesta_correctness}
  \lean{Kimchi.Commitment.IPA.ipaVesta_correctness}
  \uses{thm:ipa_correctness, def:ipa_vesta, def:commitment_perfect_correctness}
  The Vesta IPA scheme of \cref{def:ipa_vesta} satisfies perfect correctness
  (\cref{def:commitment_perfect_correctness}).
  \begin{proof}
    \uses{thm:ipa_correctness}
    Identical to the Pallas case (\cref{thm:ipa_pallas_correctness}), specializing
    the generic theorem \cref{thm:ipa_correctness} at the Vesta point group
    (\cref{def:vesta_group}) and Vesta scalar field
    (\cref{def:vesta_scalar_field}).
  \end{proof}
\end{theorem}

\begin{theorem}[evaluation binding over Vesta]
  \label{thm:ipa_vesta_binding}
  \lean{Kimchi.Commitment.IPA.ipaVesta_binding}
  \uses{thm:ipa_binding, def:ipa_vesta, def:commitment_binding}
  The Vesta IPA scheme of \cref{def:ipa_vesta} satisfies evaluation binding
  (\cref{def:commitment_binding}).
  \begin{proof}
    \uses{thm:ipa_binding}
    Identical to the Pallas case (\cref{thm:ipa_pallas_binding}), specializing the
    generic theorem \cref{thm:ipa_binding} at the Vesta point group
    (\cref{def:vesta_group}) and Vesta scalar field
    (\cref{def:vesta_scalar_field}).
  \end{proof}
\end{theorem}
